\documentclass[12pt]{article}

% --- Packages ---
\usepackage[utf8]{inputenc}
\usepackage[T1]{fontenc}
\usepackage[french]{babel}
\usepackage{amsmath, amssymb}
\usepackage{geometry}
\usepackage{tikz}
\usepackage[most]{tcolorbox}
\usepackage{multicol}
\usepackage{lmodern}
\usepackage{pgfplots}
\pgfplotsset{compat=1.18}
\usepackage{mdframed}
\usepackage{tkz-tab}
\usepackage{bm}
\usepackage{xcolor}
\usepackage{sectsty}
\usepackage[utf8]{inputenc}
\usepackage[T1]{fontenc}
\usepackage[inline]{enumitem}
\usepackage[dvipsnames]{xcolor}
\usepackage{multido}
% --- Page ---
\geometry{margin=2cm}

% Footer
\usepackage{fancyhdr}
\pagestyle{fancy}
\cfoot{\thepage}
\rfoot{\textit{G.M.}}
\renewcommand{\headrulewidth}{0pt}
\renewcommand{\footrulewidth}{0pt}

% Pour les normes : 
\newcommand*{\norme}[1]{\left\lVert\overrightarrow{#1}\right\rVert}

\usetikzlibrary{shapes.misc}

\tikzset{cross/.style={cross out, draw=black, minimum size=2*(#1-\pgflinewidth), inner sep=0pt, outer sep=0pt},
%default radius will be 1pt. 
cross/.default={4pt}}

% mettre dans le préambule, après \usepackage{xcolor}
\usepackage{sectsty} % gère facilement la couleur des titres
\setlength{\parindent}{0pt}

\sectionfont{\color{red}\Large\bfseries}     % pour \section*
\subsectionfont{\color{red}\large\bfseries\hspace{1cm}}

\begin{document}

% ============================
%   INÉQUATIONS
% ============================

\newpage

\begin{center}
    {\LARGE \textbf{Inéquations}}\\[4mm]
    {\large Seconde}
\end{center}

\hrule
\vspace{0.4cm}

% ============================
%   PARTIE 1 : Inéquations du premier degré
% ============================

\section*{\textcolor{red}{Partie 1 : Inéquations du premier degré}}

\vspace{5pt}

\begin{mdframed}[
    linewidth=1pt,
    linecolor=ForestGreen,
    innertopmargin=10pt,
    innerbottommargin=10pt
]
{\color{ForestGreen}\textbf{Définition :}}  
Une \textbf{inéquation du premier degré} est une inéquation qui peut se mettre sous la forme :
\[
ax + b \ \square \ 0
\]
où :
\begin{itemize}
    \item $a$ et $b$ sont des nombres réels avec $a \neq 0$
    \item $\square$ désigne l'un des symboles : $<$, $\le$, $>$ ou $\ge$
\end{itemize}
\end{mdframed}

\vspace{10pt}

\underline{Exemple :}
\[
3x - 7 \ge 5
\]

\vspace{5pt}

\underline{Résolution :}
\[
3x - 7 \ge 5
\]
\[
3x \ge 12
\]
\[
x \ge 4
\]

\vspace{10pt}

\begin{mdframed}[
    linewidth=1pt,
    linecolor=Red,
    innertopmargin=10pt,
    innerbottommargin=10pt
]
{\color{red}
\underline{Attention :}
Lorsque l'on \textbf{multiplie ou divise} une inéquation par un nombre \textbf{négatif},  
le \textbf{signe de l'inégalité change de sens}.
}
\end{mdframed}

\vspace{5pt}

\begin{mdframed}[
    linewidth=1pt,
    linecolor=RoyalBlue,
    innertopmargin=10pt,
    innerbottommargin=10pt
]
{\color{RoyalBlue}\textbf{\underline{Exercice d'application n°1 :}}}
\vspace{5pt}

Résoudre les inéquations suivantes et représenter les solutions sur une droite graduée :
\begin{enumerate}
    \item $2x + 3 < 4 - 5x$
    \item $2(x - 4) \leq 4x - 5$
\end{enumerate}
\end{mdframed}

\vspace{10pt}
\newpage
% ============================
%   PARTIE 2 : Tableau de signes
% ============================

\section*{\textcolor{red}{Partie 2 : Tableau de signes}}

\vspace{5pt}

\subsection*{2.1 Expression de la forme $ax+b$}

\underline{Exemple d'introduction :}

\vspace{5pt}

On souhaite déterminer le signe de l'expression :
\[
2x - 10
\]

\vspace{5pt}

\underline{Étape 1 :} Déterminons pour quelle valeur de $x$ l'expression $2x - 10$ s'annule.
\[
\text{On résout }2x - 10 = 0
\]
\vspace{-20pt}
{\color{blue}
\[
2x = 10
\]
\[
x = \dfrac{10}{2}
\]
\[
x = 5
\]
}
% \\$2x = 10$
%     \\$x = \dfrac{10}{2}$
%     \\$x = 5$
\vspace{5pt}

\underline{Étape 2 :} On dresse le tableau de signes de $2x - 10$
\vspace{5pt}\\
Sachant que $x \longmapsto 2x - 10$ est une fonction affine représentée par 
\textbf{une droite}, 
on peut 
déduire le tableau de signe suivant : 
\[
\begin{tikzpicture}
   \tkzTabInit{$x$ / 1 , $2x-10$ / 1}{$-\infty$, ${\color{blue}5}$, $+\infty$}
   \tkzTabLine{, {\color{red}-}, z, {\color{red}+}, }
\end{tikzpicture}
\]

\vspace{5pt}

En traçant le graphique de $x \longmapsto 2x - 10$, on retrouve ce 
résultat.

\begin{mdframed}[
    linewidth=1pt,
    roundcorner=5pt,
    linecolor=red
]
    {\color{red}\underline{À RETENIR :}}
    Pour les fonctions de la forme $ax+b$ :
    \begin{itemize}
        \item si $\color{red}a>0$, alors on commence avec un $\color{red}-$ dans le 
        tableau de signe, 
        puis un $\color{red}+$
        \item si $\color{red}a<0$, alors on commence avec un $\color{red}+$ dans le 
        tableau de signe, puis un $\color{red}-$
    \end{itemize}
\end{mdframed}

\begin{center}
\begin{tikzpicture}
\begin{axis}[
    axis lines=middle,       % axes au milieu
    xlabel=$x$, ylabel=$y$,
    xmin=-1, xmax=7,         % bornes x
    ymin=-4, ymax=4,         % bornes y
    width=12cm,               % largeur du graphique
    height=6cm,               % hauteur du graphique
]

% tracer la droite
\addplot[thick,red,domain=-1:7,samples=2]{2*x - 10};

% marquer la racine x=5
\addplot[blue,mark=*] coordinates {(5,0)};
\addplot[blue,dashed] coordinates {(5,-4) (5,4)};

\end{axis}
\end{tikzpicture}
\end{center}


\vspace{10pt}


\begin{mdframed}[
    linewidth=1pt,
    linecolor=RoyalBlue,
    innertopmargin=10pt,
    innerbottommargin=10pt
]
{\color{RoyalBlue}\textbf{\underline{Exercice n°2 :}}}
\vspace{5pt}

Déterminer le signe d'une expression de la forme $ax+b$
\vspace{-10pt}
\begin{multicols}{2}
    \begin{enumerate}
        \item $2x + 6$
        \item $-3x + 12$
    \end{enumerate}
\end{multicols}

\end{mdframed}

\vspace{10pt}

\subsection*{\textcolor{red}{2.2 Expression du type $(ax+b)(cx+d)$}}

\vspace{5pt}

\underline{Exemple :}
\[
(2x - 1)(x + 3)
\]

\vspace{5pt}

\underline{But : }
On cherche le signe de l'expression $(2x - 1)(x + 3)$, c'est-à-dire \textbf{pour quelles
valeurs de x} elle est positive ou négative.
\vspace{5pt}\\
Le signe de $(2x - 1)(x + 3)$ dépend \_\_\_\_\_\_\_\_\_\_\_\_\_\_\_\_\_\_\_\_\_\_\_\_\_\_\_\_\_\_\_\_\_\_\_\_

\vspace{15pt}
\underline{On résout :}
\begin{minipage}[t]{.35\textwidth}%
% \vspace(0pt)
    \[
    \begin{aligned}
        2x - 1 &= 0 \\
        2x &= 1 \\
        x &= \dfrac{1}{2}
    \end{aligned}
    \]
\end{minipage}%
\hfill
\begin{minipage}[t]{.55\textwidth}%
% \vspace(0pt)
    \[
    \begin{aligned}
        x + 3 &= 0 \\
        x &= -3
    \end{aligned}
    \]
\end{minipage}%

\vspace{15pt}

\underline{On peut 
déduire le tableau de signe suivant :}
\[
\begin{tikzpicture}
   \tkzTabInit[lgt = 3.5]{$x$ / 1 , $2x - 1$ / 1, $x + 3$ / 1, $(2x-1)(x+3)$ / 1}
   {$-\infty$, $-3$, $\dfrac{1}{2}$, $+\infty$}
   \tkzTabLine{, {\color{red}-}, , {\color{red}-}, z, {\color{red}+}, }
   \tkzTabLine{, {\color{red}-}, z, {\color{red}+}, , {\color{red}+}, }
   \tkzTabLine{, {\color{red}+}, z, {\color{red}-}, z, {\color{red}+}, }
\end{tikzpicture}
\]

\begin{mdframed}[
    linewidth=1pt,
    roundcorner=5pt,
    linecolor=red,
    innerbottommargin=10pt,
    innertopmargin=10pt
]
    {\color{red}\underline{À RETENIR :}}
    Pour la dernière ligne du tableau, on utilise la règle des signes !
\end{mdframed}

\begin{mdframed}[
    linewidth=1pt,
    linecolor=RoyalBlue,
    innertopmargin=10pt,
    innerbottommargin=10pt
]
{\color{RoyalBlue}\textbf{\underline{Exercice n°3 :}}}
\vspace{5pt}

Dresser le tableau de signes de l'expression $(3x-9)(1-2x)$.
\end{mdframed}

\vspace{20pt}
\newpage

% ============================
%   PARTIE 3 : Inéquation produit
% ============================

\section*{\textcolor{red}{Partie 3 : Inéquation produit}}

\vspace{5pt}

\underline{Exemple :}
\vspace{5pt}\\On veut résoudre :
\begin{center}
    $(x-3)(1-3x) > 0$
\end{center}

À nouveau, le signe du produit dépend du signe de chaque facteur.

\vspace{10pt}
\underline{On résout :}\\
\begin{minipage}[t]{.25\textwidth}%
% \vspace(0pt)
    \[
    \begin{aligned}
        x-3 &= 0\\
        x &= 3
    \end{aligned}
    \]
\end{minipage}%
\hfill
\begin{minipage}[t]{.75\textwidth}%
% \vspace(0pt)
    \[
    \begin{aligned}
        1-3x &= 0\\
        x &= \dfrac{1}{3}
    \end{aligned}
    \]
\end{minipage}%

\vspace{15pt}

\underline{On peut 
déduire le tableau de signe suivant :}
\[
\begin{tikzpicture}
   \tkzTabInit[lgt = 3.5]{$x$ / 1 , $x-3$ / 1, $1 - 3x$ / 1, $(x-3)(1-3x)$ / 1}
   {$-\infty$, $\dfrac{1}{3}$, $3$, $+\infty$}
   \tkzTabLine{, {\color{blue}-}, , {\color{blue}-}, z, {\color{blue}+}, }
   \tkzTabLine{, {\color{blue}+}, z, {\color{blue}-}, , {\color{blue}-}, }
   \tkzTabLine{, {\color{blue}-}, z, {\color{red}+}, z, {\color{blue}-}, }
\end{tikzpicture}
\]

On lit facilement dans le tableau que $(x-3)(1-3x) > 0$ lorsque $\dfrac{1}{3} < x < 3$.
\vspace{5pt}\\L'ensemble des solutions de l'inéquation est donc :
\[
S = \left]\dfrac{1}{3};3\right[
\]


\begin{mdframed}[
    linewidth=1pt,
    linecolor=RoyalBlue,
    innertopmargin=10pt,
    innerbottommargin=10pt
]
{\color{RoyalBlue}\textbf{\underline{Exercice n°4 :}}}
\vspace{5pt}

Résoudre dans $\mathbb{R}$ l'inéquation suivante :
\[
(3 - 6x)(x+2) \geq 0
\]
\end{mdframed}

\vspace{10pt}

% ============================
%   PARTIE 4 : Inéquation quotient
% ============================

\section*{\textcolor{red}{Partie 4 : Inéquation quotient}}

\vspace{5pt}

\underline{Exemple :}
\vspace{5pt}\\On veut résoudre :
\begin{center}
    $\dfrac{2x - 3}{1 - 4x} \leq 0$
\end{center}

Le signe de $\dfrac{2x - 3}{1 - 4x}$ dépend du signe des expressions $2x - 3$ et $1 - 4x$.

\vspace{5pt}{\color{red}\textbf{ATTENTION} : le dénominateur doit être différent de $0$ !}

\vspace{10pt}
\underline{On résout :}\\
\begin{minipage}[t]{.25\textwidth}%
% \vspace(0pt)
    \[
    \begin{aligned}
        2x-3 &= 0\\
        x &= \dfrac{3}{2}
    \end{aligned}
    \]
\end{minipage}%
\hfill
\begin{minipage}[t]{.75\textwidth}%
% \vspace(0pt)
    \[
    \begin{aligned}
        1-4x &\neq 0\\
        x &\neq \dfrac{1}{4}
    \end{aligned}
    \]
\end{minipage}%

\vspace{15pt}

\underline{On peut 
déduire le tableau de signe suivant :}
\[
\begin{tikzpicture}
    \tkzTabInit[lgt = 3.5]{$x$ / 1 , $2x-3$ / 1, $1 - 4x$ / 1, 
        $\dfrac{2x-3}{1-4x}$ / 1.5}
    {$-\infty$, $\dfrac{1}{4}$, $\dfrac{3}{2}$, $+\infty$}
    \tkzTabLine{, {\color{blue}-}, , {\color{blue}-}, z, {\color{blue}+}, }
    \tkzTabLine{, {\color{blue}+}, z, {\color{blue}-}, , {\color{blue}-}, }
    \tkzTabLine{, {\color{red}-}, d, {\color{blue}+}, z, {\color{red}-}, }
\end{tikzpicture}
\]

{\color{red}\textbf{\underline{À RETENIR :}} dans la dernière ligne, 
$\dfrac{1}{4}$ est une 
\textbf{valeur interdite} ! On ne peut pas diviser par $0$! On place donc deux 
traits parallèles dans le tableau de 
signes.}

\vspace{10pt}

On lit facilement dans le tableau que $\dfrac{2x - 3}{1 - 4x} \leq 0$ 
lorsque $x < \dfrac{1}{4}$ \textbf{ET} $x \geq \dfrac{3}{2}$.
\vspace{5pt}\\L'ensemble des solutions de l'inéquation est donc :
\[
S = \left]-\infty;\dfrac{1}{4}\right[ \bigcup \left[\dfrac{3}{2};+\infty\right[
\]


\begin{mdframed}[
    linewidth=1pt,
    linecolor=RoyalBlue,
    innertopmargin=10pt,
    innerbottommargin=10pt
]
{\color{RoyalBlue}\textbf{\underline{Exercice n°4 :}}}
\vspace{5pt}

Résoudre dans $\mathbb{R}$ l'inéquation suivante :
\[
\dfrac{2-6x}{3x-2} < 0
\]
\end{mdframed}


\end{document}
